\chapter{Cleft palate}

\section*{Definition and Overview}
A cleft palate is a congenital craniofacial anomaly characterised by an opening or split in the roof of the mouth (the palate). This condition occurs during embryonic development when the tissue that forms the hard and soft palates fails to fuse properly. A cleft palate can involve the soft palate (the muscular back part of the mouth), the hard palate (the bony front part), or both, and may extend forward into the alveolar ridge (gum). It can present in isolation (isolated cleft palate) or in association with a cleft lip. A cleft palate significantly impacts physiological functions, particularly the ability to feed, speak, hear, and develop normal dentition. Effective management requires a lifelong, coordinated multidisciplinary team approach starting at birth.

\section*{Epidemiology}
Isolated cleft palate occurs in approximately 1 in 1,500 to 1 in 2,000 live births globally. Unlike cleft lip (which is more common in males), isolated cleft palate is more common in females, with a female-to-male ratio of approximately 2:1. In Australia, the incidence of cleft palate remains consistent with international averages. Isolated cleft palate has a much higher association with genetic syndromes and other congenital anomalies than cleft lip; up to 50\% of cleft palate cases are syndromic. Common associated conditions include Pierre Robin sequence, Stickler syndrome, velocardiofacial syndrome (22q11.2 deletion syndrome), and Treacher Collins syndrome. The remaining cases are non-syndromic and arise from multifactorial genetic and environmental interactions.

\section*{Aetiology and Pathophysiology}
The development of a cleft palate is a complex process involving the failure of palatal shelf elevation and fusion during early foetal growth.
\begin{itemize}
    \item \textbf{Embryological fusion failure:} Pathophysiologically, the palate forms between the eighth and twelfth weeks of gestation. A cleft palate results when the lateral palatine processes (palatal shelves) fail to elevate, meet, and fuse with each other, the nasal septum, or the primary palate.
    \item \textbf{Genetic syndromes and mutations:} Mutations in genes such as $TBX22$ (X-linked cleft palate), $IRF6$, $GRHL3$, and chromosomal microdeletions (e.g., 22q11.2) are major genetic drivers.
    \item \textbf{Maternal risk factors:} Environmental factors including maternal smoking, heavy alcohol consumption, and uncontrolled gestational diabetes significantly increase the risk of palatal clefting.
    \item \textbf{Nutritional deficiencies:} Inadequate maternal intake of folic acid (vitamin B9) and other essential micronutrients during the first trimester impairs cellular proliferation and fusion.
    \item \textbf{Teratogenic drug exposure:} First-trimester maternal exposure to medications such as benzodiazepines, corticosteroids, and certain anticonvulsants (e.g., topiramate) is linked to a higher risk of cleft palate.
    \item \textbf{Mechanical interference:} In conditions like Pierre Robin sequence, a small lower jaw (micrognathia) prevents the tongue from descending, which mechanically blocks the palatal shelves from fusing.
\end{itemize}

\section*{Clinical Features}
A cleft palate may not be immediately visible externally but is identified upon physical inspection of the oral cavity at birth or via prenatal imaging.
\begin{itemize}
    \item \textbf{Split in the roof of the mouth:} A visible and palpable gap in the hard and/or soft palate, which may be unilateral or bilateral.
    \item \textbf{Submucous cleft palate:} A specific variant where the overlying oral mucosa is intact, but the underlying palatal muscles are cleft, often characterised by a bifid uvula, a translucent zone in the soft palate (zona pellucida), and a notch in the hard palate.
    \item \textbf{Severe feeding difficulties:} Inability to create the negative intraoral pressure (suction) necessary for breastfeeding or normal bottle-feeding, leading to nasal regurgitation of milk, prolonged feeding times, fatigue, and potential choking.
    \item \textbf{Speech anomalies:} Development of hypernasal speech and compensatory articulation errors due to the inability of the soft palate to seal the oral cavity from the nasal cavity during speech (velopharyngeal insufficiency).
    \item \textbf{Otological issues:} High incidence of middle ear fluid accumulation (otitis media with effusion) due to dysfunction of the tensor veli palatini muscle, which controls the Eustachian tube.
\end{itemize}

\section*{Differential Diagnosis}
Clinicians must evaluate infants with a cleft palate for associated systemic features to differentiate isolated cases from genetic syndromes.
\begin{itemize}
    \item \textbf{Isolated cleft palate:} A cleft palate present without any other associated physical, intellectual, or developmental anomalies.
    \item \textbf{Pierre Robin sequence:} A triad consisting of micrognathia (small lower jaw), glossoptosis (tongue falling back into the throat), and airway obstruction, frequently presenting with a U-shaped cleft palate.
    \item \textbf{Velocardiofacial syndrome:} A genetic syndrome presenting with cleft palate (or submucous cleft), congenital heart defects (e.g., Tetralogy of Fallot), learning difficulties, immune deficiency, and characteristic facial features.
    \item \textbf{Stickler syndrome:} A genetic connective tissue disorder characterized by cleft palate, severe myopia (nearsightedness), retinal detachment, hearing loss, and early-onset joint problems.
    \item \textbf{Congenital palatal insufficiency:} A condition where the palate is intact but structurally short or functionally weak, causing hypernasal speech similar to a cleft palate.
\end{itemize}

\section*{When to Seek Medical Care}
Immediate medical consultation is required after birth if a cleft palate is detected or suspected (e.g., if a baby has a split uvula or milk regularly comes out of their nose during feeding).

\textbf{Call 000 (or local emergency services) for:}
\begin{itemize}
    \item Severe respiratory distress, choking, or gasping for air, which is a particular risk in infants with Pierre Robin sequence.
    \item Blue or pale discolouration of the skin, lips, or tongue (cyanosis).
    \item Persistent choking or aspiration of milk leading to coughing, wheezing, or difficulty breathing.
    \item Signs of severe systemic infection or lethargy in the infant.
\end{itemize}

\section*{Diagnosis and Investigations}
Diagnosis is confirmed through clinical examination and developmental assessments.
\begin{itemize}
    \item \textbf{Neonatal oral examination:} Visual inspection and digital palpation of the entire roof of the mouth immediately after birth, which is essential to detect posterior or submucous clefts.
    \item \textbf{Feeding evaluation:} Assessment by a speech pathologist or lactation consultant to evaluate sucking mechanism, airway safety, and to prescribe specialized feeding equipment.
    \item \textbf{Audiological assessment:} Baseline and serial hearing tests (including tympanometry and brainstem auditory response) to monitor for middle ear fluid and conductive hearing loss.
    \item \textbf{Genetic screening and microarray:} Highly recommended for all cases of isolated cleft palate to screen for syndromic etiologies, especially 22q11.2 microdeletions.
    \item \textbf{Speech and language evaluation:} Routine assessment beginning around 12 months of age to monitor speech milestones and identify velopharyngeal insufficiency.
\end{itemize}

\section*{Management}
Management is multidisciplinary and involves sequential interventions from infancy to adulthood.
\begin{itemize}
    \item \textbf{Specialist feeding support:} Use of specialized squeeze bottles (such as the SpecialNeeds Feeder) and orthodontic palatal obturators to facilitate feeding without requiring suction.
    \item \textbf{Primary surgical repair (Palatoplasty):} Surgical closure of the cleft palate, typically performed between 9 and 18 months of age, before the child develops significant speech patterns, to optimize velopharyngeal function.
    \item \textbf{Tympanostomy tubes (Grommets):} Surgical insertion of small tubes into the eardrum to drain middle ear fluid and prevent recurrent infections and hearing loss, often performed at the time of palatoplasty.
    \item \textbf{Speech pathology intervention:} Intensive speech therapy to correct compensatory articulation habits and manage hypernasality.
    \item \textbf{Orthodontic and dental care:} Comprehensive monitoring of dental arches, bone grafting of the alveolar ridge (if affected, usually between 8 and 11 years of age), and braces.
    \item \textbf{Pharyngeal flap surgery or sphincter pharyngoplasty:} Secondary surgical procedures performed in childhood if velopharyngeal insufficiency persists after primary repair.
\end{itemize}

\section*{Complications}
\begin{itemize}
    \item \textbf{Velopharyngeal insufficiency:} Inability of the repaired soft palate to seal against the back of the throat, resulting in hypernasal speech and air escape through the nose.
    \item \textbf{Conductive hearing loss:} Persistent fluid in the middle ear (glue ear) due to Eustachian tube dysfunction, which can impair language acquisition.
    \item \textbf{Malnutrition and failure to thrive:} Severe early feeding difficulties leading to inadequate caloric intake and poor growth.
    \item \textbf{Dental anomalies and malocclusion:} Missing, extra, or poorly positioned teeth and underdeveloped upper jaw (maxillary hypoplasia), requiring complex orthodontic treatment and jaw surgery (osteotomy).
    \item \textbf{Psychosocial and educational challenges:} Social integration difficulties, communication barriers, and reduced self-esteem.
\end{itemize}

\section*{Prognosis and Outcomes}
With modern multidisciplinary care, the prognosis for children born with a cleft palate is generally excellent. The majority of children achieve normal or near-normal speech, hearing, and dental function by school age. Early surgical intervention (before 18 months) is critical for optimal speech outcomes, with only 10--20\% of children requiring secondary surgery for speech correction. While orthodontic treatment and dental monitoring are often required throughout childhood and adolescence, long-term functional and aesthetic outcomes are highly successful.

\section*{Prevention and Self-Care}
\begin{itemize}
    \item \textbf{Preconceptional folic acid:} Women of childbearing age should take 400\,microg of folic acid daily starting at least one month before conception to reduce the risk of neural tube and craniofacial defects.
    \item \textbf{Avoid known teratogens:} Eliminate smoking, alcohol consumption, and non-prescribed medications during pregnancy, especially in the first trimester.
    \item \textbf{Obtain genetic counselling:} Families with a history of clefting or known genetic syndromes should seek genetic counselling to understand recurrence risks.
    \item \textbf{Learn specialized feeding techniques:} Parents should be trained in upright feeding positions, frequent burping, and using specialized squeeze bottles.
    \item \textbf{Promote speech development:} Engage in regular talking, reading, and singing with your child, and follow home-based speech therapy activities recommended by your therapist.
\end{itemize}

\section*{Sources and Further Reading}
The following reputable organisations provide further information on cleft palate:
\begin{itemize}
    \item CleftPALS (Cleft Palate and Lip Society Australia). \textit{Information, feeding equipment, and family support services.} https://www.cleftpals.org.au/
    \item Sydney Children's Hospitals Network. \textit{Comprehensive clinical resource on cleft palate management.} https://www.schn.health.nsw.gov.au/
    \item Royal Children's Hospital Melbourne. \textit{Cleft palate repair clinical guidelines and parent resources.} https://www.rch.org.au/
    \item Cleft Lip and Palate Association (CLAPA). \textit{International resources, feeding guides, and patient advocacy.} https://www.clapa.com/
    \item Mayo Clinic. \textit{Patient education on cleft palate symptoms, diagnosis, and treatment.} https://www.mayoclinic.org/
\end{itemize}
